\documentclass{article}
\usepackage[utf8]{inputenc}
\usepackage{array}
\newcolumntype{C}[1]{>{\centering\arraybackslash}m{#1}}
\usepackage{amsmath}
\usepackage{relsize}
\usepackage{graphicx}

\usepackage[noabbrev]{cleveref}

\crefformat{section}{\S#2#1#3}
\crefformat{subsection}{\S#2#1#3}
\crefformat{subsubsection}{\S#2#1#3}
\crefformat{figure}{#2Figura~#1#3}
\crefformat{footnote}{#2\footnotemark[#1]#3}
\crefformat{table}{#2Tabella~#1#3}
\crefformat{equation}{#2Equazione~#1#3}
\crefformat{chapter}{#2Capitolo~#1#3}

\usepackage{tikz}
\usetikzlibrary{trees}

\begin{document}


\include{title}
\thispagestyle{empty}
\newpage

\section{Problema 1}
\paragraph{Traccia} Si consideri il seguente problema di minimo percorso vincolato. Sia data una matrice
rettangolare di 1 e 0, in cui 1 indica una cella che è possibile percorrere e 0 una cella che
non è possibile percorrere. Inoltre, si assuma che non è possibile percorrere neanche le
quattro celle adiacenti (sinistra, destra, sopra e sotto) ad una cella che contiene uno 0.
Calcolare la lunghezza del più breve percorso possibile da qualsiasi cella nella prima
colonna a qualsiasi cella nell'ultima colonna della matrice. L’obiettivo è evitare le celle
contrassegnate con 0 e le loro quattro celle adiacenti (sinistra, destra, sopra e sotto).
Le mosse possibili da una determinata cella sono lo spostamento di una cella adiacente a
sinistra, destra, sopra o sotto (non sono consentiti spostamenti in diagonale).

\paragraph{Soluzione}

La soluzione prevede l'utilizzo della tecnica del \textit{backtracking}. Viene costruita una matrice di appoggio \textit{visited} i cui elementi sono $0$ se la corrispondente cella della matrice di ingresso non è stata ancora visitata, o $1$ altrimenti. Viene richiamata la funzione \textit{find\_min\_path} ricorsivamente, testando ogni volta una delle quattro direzioni con la funzione \textit{is\_valid}. Se la cella scelta soddisfa il vincolo posto dal problema (le celle in torno agli $0$ non sono valide e non si violano i confini della matrice), allora l'algoritmo prosegue in quella direzione, altrimenti viene annullato l'ultimo movimento effettuato e si prova un'altra direzione.
La funzione si ferma nel momento in cui si arriva alla cella di destinazione con coordinate \textit{(end\_x, end\_y)} e aggiorna la variabile risultato quando la lunghezza del percorso è minore del risultato precedente (la variabile risultato viene inizializzata a \textit{INT\_MAX}).
La funzione \textit{solution} trova tutte le soluzioni combinando tutti gli elementi della prima colonna con quelli dell'ultima come punto di partenza e di arrivo, stampando il risultato se esiste, o un messaggio di errore nel caso in cui la variabile risultato sia ancora uguale a \textit{INT\_MAX}.

\paragraph{Complessità}
La complessità dell'algoritmo dipende da due fattori: il primo è la funzione di \textit{backtracking} \textit{find\_min\_path }, che genera la ricorrenza $T(P) = 4\cdot T(P-1) + O(1)$ dove $P$ è uguale al numero di celle della matrice ($NxM$). Questa funzione verrà chiamata per ogni combinazione degli elementi della prima colonna e quelli dell'ultima, cioè $2^M$ volte. In totale quindi si avrà una complessità pari a $O(2^M \cdot 4^{NxM})$

\newpage
\section{Problema 2}
\paragraph{Traccia}
Sia dato un insieme di n elementi, e si supponga che ogni elemento possa essere
accoppiato con qualche altro elemento oppure può essere non accoppiato (“single”). Ogni
elemento può essere accoppiato solo una volta. Si implementi un algoritmo per scoprire il
numero totale di modi in cui gli n elementi possono restare single o possono essere
accoppiati. Si utilizzi la programmazione dinamica per risolvere il problema.

\paragraph{Soluzione}
La soluzione prevede l'utilizzo della programmazione dinamica come richiesto dalla traccia. Il problema inizialmente è stato risolto con la semplice ricorsione: ogni partecipante può o stare da solo, quindi il problema si riduce di un elemento, oppure può appaiarsi con un altro degli $n-1$ partecipanti: in tal caso si avranno $n-1$ combinazioni della prima coppia moltiplicato con il problema iniziale diminuito di due elementi. In definitiva avremo quindi la ricorrenza: $T(n) = T(n-1) + (n-1) \cdot T(n-2)$. Analogamente al caso studiato per la serie di Fibonacci, si può applicare la programmazione dinamica in modo da conservare gli elementi già calcolati in un vettore \textit{pairs} e sostituire alle chiamate ricorsive la lettura dal \textit{pairs} esclusi i casi $n=0,\;n=1$.

\paragraph{Complessità}
Con l'applicazione della programmazione dinamica, il problema si risolve con un unico ciclo sugli $n$ elementi, dove $n$ è il parametro di ingresso della soluzione, ottenendo quindi una complessità $O(n)$.

\newpage
\section{Problema 3}
Si consideri una matrice di 0 ed 1, in cui “1” indica “posizione occupata” e “0” indica
“posizione libera”. Si scriva un algoritmo per determinare la sotto-matrice massima libera
(ossia che contiene tutti 0). L’algoritmo deve riportare il numero di 0 di tale sotto-matrice.

\paragraph{Soluzione}
La soluzione prevede l'utilizzo della programmazione dinamica. In particolare, per ogni riga, partendo dall'ultima colonna a destra, viene contato il numero di zeri contigui e il risultato viene memorizzato in una matrice di appoggio \textit{dp}. In seguito, la matrice \textit{dp} viene controllata per righe da sinistra verso destra e si calcolano due variabili: \textit{columns} e \textit{area}. \textit{Columns} è inizializzata al numero di colonne della matrice di partenza e ad ogni iterazione viene scelta confrontando il valore precedente e il valore corrente della matrice \textit{dp} scegliendo il minimo,come illustrato in \cref{fig:notebook}. \textit{Area} viene calcolata moltiplicando \textit{columns} e il valore della riga corrente e corrisponde al numero di zeri nella sotto-matrice.

\begin{figure}[h]
    \centering
    \includegraphics[width=14cm]{amogus-matrix.png}
    \caption{Schema della soluzione}
    \label{fig:notebook}
\end{figure}

\paragraph{Complessità}
Nella prima parte della soluzione si è popolata la matrice di appoggio \textit{dp} con una complessità $O(mn)$. Successivamente, vengono eseguiti 3 cicli innestati di cui i primi due sono analoghi ai precedenti, mentre l'ultimo esegue al più $m$ volte essendo il caso peggiore quello in cui la matrice sia composta da soli zeri. In definitiva k, la complessità totale risulta essere $O(m^2n)$ invece della soluzione \textit{bruteforce} che prevede una complessità pari a $O((mn)^3)$.

\end{document}
